\documentclass[11pt]{article}
\usepackage[T1]{fontenc}
\usepackage[utf8]{inputenc}
\usepackage{lmodern}
\usepackage[margin=0.85in]{geometry}
\usepackage{amsmath,amssymb,booktabs,array,tabularx}
\usepackage{microtype}
\usepackage[hidelinks]{hyperref}
\usepackage{caption}
\captionsetup{font=small,labelfont=bf}
\setlength{\parindent}{0pt}
\setlength{\parskip}{5pt}
\setlength{\emergencystretch}{2em}
\newcommand{\ex}{\mathrm{e}}
\begin{document}
\begin{center}
    {\Large\bfseries Problem Set 8}\\[4pt]
    Questions (d) and (e)
\end{center}

\section*{Question (d)}

Each month, run the cross-sectional regression across the 75 portfolios:
\begin{equation}
 r_{i,t}-r_{f,t}=a_t+b_{1,t}X_{1,i,t}+b_{2,t}X_{2,i,t}
                    +b_{3,t}X_{3,i,t}+\varepsilon_{i,t},
 \label{eq:fmb}
\end{equation}
where returns are in decimal units and
\begin{equation}
\begin{aligned}
 X_{1,i,t}&=r_{i,t-1},\\
 X_{2,i,t}&=\prod\nolimits_{j=2}^{12}(1+r_{i,t-j})-1,\\
 X_{3,i,t}&=\prod\nolimits_{j=13}^{60}(1+r_{i,t-j})-1.
\end{aligned}
\label{eq:predictors}
\end{equation}
The second and third predictors compound 11 and 48 monthly returns, respectively.
Missing-value codes $-999$ and $-99.99$ are excluded before constructing the
predictors. Retaining months with complete observations for all 75 portfolios
gives 1,000 monthly regressions, from March 1943 through June 2026.
For each coefficient, report its time-series mean and conventional
Fama--MacBeth $t$-statistic:
\begin{equation}
 \bar b_j=\frac{1}{T}\sum_{t=1}^{T}b_{j,t},\qquad
 t(\bar b_j)=\frac{\bar b_j}{s(b_{j,t})/\sqrt{T}},\qquad T=1{,}000,
 \label{eq:fmb-stat}
\end{equation}
where $s$ is the sample standard deviation; the same calculation applies to the intercept.

\begin{table}[htbp]
\centering
\caption{Time-series averages of monthly cross-sectional coefficients}
\label{tab:fmb}
\begin{tabular}{lrr}
\toprule
Variable & Average coefficient & Fama--MacBeth $t$-statistic\\
\midrule
Intercept & 0.003070 & 2.0858\\
Past one-month return & 0.053647 & 4.8898\\
Past return, $t-12$ to $t-2$ & 0.018197 & 6.4683\\
Past return, $t-60$ to $t-13$ & 0.002779 & 3.6992\\
\bottomrule
\end{tabular}
\end{table}

The Fama--MacBeth coefficients are positive and statistically significant at all
three horizons. The past one-month coefficient is 0.05365 ($t=4.89$), indicating
short-term continuation rather than reversal. The past 12-month coefficient,
skipping the latest month, is 0.01820 ($t=6.47$), supporting intermediate-term
momentum. The past 60-month coefficient, skipping the latest year, is 0.002779
($t=3.70$), indicating long-term continuation rather than reversal. Holding other
predictors constant, a 10-percentage-point increase in these measures predicts
approximately 53.6, 18.2, and 2.8 basis points higher monthly excess returns,
respectively. The intercept is significant at 5\% ($t=2.09$).

Overall, past winners tend to outperform past losers. Investor underreaction and
gradual information diffusion could explain this persistence; the negative
coefficients expected from overreaction followed by reversal are absent.

These results potentially challenge weak-form efficiency because historical
returns predict future returns. However, predictability does not establish
exploitable abnormal returns. The regressions subtract the risk-free rate but do
not control for systematic risk, so the patterns could reflect risk compensation.
Transaction costs, turnover, and short-selling costs could absorb trading profits.
Data mining could also exaggerate significance, making risk-adjusted,
out-of-sample results after costs necessary for a stronger conclusion.

\clearpage
\section*{Question (e)}

\subsection*{Sample splitting and strategy construction}

Split: 700 training / 300 testing months (July 2001 to June 2026); the first
120 training months calibrate the model and are excluded from performance.
The training sample runs from March 1943 through June 2001. Training performance
therefore covers March 1953 through June 2001 (580 months).

The 75 portfolios comprise the 1-month, 12-month, and 60-month sorting families.
Within each family, $S1$--$S5$ denote smallest--largest size groups, and
$P1$--$P5$ denote lowest--highest past-return groups. Use the three lagged-return
signals in Equation~\eqref{eq:predictors}.

The training walk-forward re-estimates coefficients every 12 months using prior
data only, whereas the test uses coefficients fixed at the end of training; the
two are not like-for-like. At each estimation date, retain a Fama--MacBeth slope
only if its mean is positive and its conventional $t$-statistic exceeds 1.96;
otherwise set it to zero. The momentum forecast is
\begin{equation}
 m_{i,t}=\widehat b_1X_{1,i,t}+\widehat b_2X_{2,i,t}
                              +\widehat b_3X_{3,i,t}.
 \label{eq:momentum}
\end{equation}

Seasonal increments are estimated net of the fitted momentum component.
For each portfolio, regress training excess returns minus the momentum forecast
on an intercept, a January dummy, and a December dummy; February--November is
the baseline. Retain a positive January increment with an HC1 $t$-statistic
above 1.96 only for $S1P1$ and $S1P2$ across the three families. Retain a positive
December increment with an HC1 $t$-statistic above 1.96 only for $P4$ and $P5$
in the 12-month family. Set other increments to zero.

Parameter selection optimizes the walk-forward training Sharpe. The candidate
seasonal blends are $0.5$, $1.0$, and $2.0$, and the candidate numbers of
portfolios per leg are 7 and 15. Select the highest net monthly training Sharpe;
ties favor the lower blend, then fewer portfolios per leg. The selected blend is
\textbf{2.0}, with \textbf{7 portfolios per leg}.

\subsection*{Frozen rule}

The final training slopes are
\begin{equation}
 (\widehat b_1,\widehat b_2,\widehat b_3)
 =(0.072603246787,\;0.021941645673,\;0.004497895106).
 \label{eq:frozen-slopes}
\end{equation}
Table~\ref{tab:seasonal} reports the nonzero frozen seasonal increments.
All coefficients are fixed throughout testing.

\begin{table}[htbp]
\centering
\caption{Nonzero frozen momentum-adjusted seasonal increments}
\label{tab:seasonal}
\begin{tabular}{llrr}
\toprule
Sorting family & Portfolio & January & December\\
\midrule
1-month  & $S1P1$ & 0.084157 & 0\\
1-month  & $S1P2$ & 0.056753 & 0\\
12-month & $S1P1$ & 0.089833 & 0\\
12-month & $S1P2$ & 0.061516 & 0\\
60-month & $S1P1$ & 0.084715 & 0\\
60-month & $S1P2$ & 0.054828 & 0\\
\midrule
12-month & $S1P4$ & 0 & 0.017001\\
12-month & $S1P5$ & 0 & 0.022356\\
12-month & $S2P4$ & 0 & 0.022110\\
12-month & $S2P5$ & 0 & 0.025407\\
12-month & $S3P4$ & 0 & 0.020930\\
12-month & $S3P5$ & 0 & 0.024924\\
12-month & $S4P4$ & 0 & 0.019643\\
12-month & $S4P5$ & 0 & 0.022779\\
12-month & $S5P4$ & 0 & 0.013961\\
12-month & $S5P5$ & 0 & 0.019163\\
\bottomrule
\end{tabular}
\par\smallskip
\begin{minipage}{0.9\textwidth}
\footnotesize
Increments are in decimal monthly-return units, displayed to six decimal
places. All unlisted increments are zero. The strategy uses the unrounded
training estimates.
\end{minipage}
\end{table}

Let $d_{i,t}$ equal the retained January increment in January, the retained
December increment in December, and zero otherwise. The monthly score is
\begin{equation}
 q_{i,t}=z(m_{i,t})+2z(d_{i,t}),\qquad
 z(v_i)=\frac{v_i-\bar v}{\sqrt{75^{-1}\sum_{j=1}^{75}(v_j-\bar v)^2}}.
 \label{eq:score}
\end{equation}
Set a component to zero if its cross-sectional standard deviation is no greater
than $10^{-12}$. Long the seven highest scores at $+1/7$ each and short the
seven lowest scores at $-1/7$ each; other risky weights are zero. Shorts are
relative hedges. If the score range is no greater than $10^{-12}$, hold cash.
For ties, use a stable ascending sort in portfolio order: 1-month, 12-month,
60-month families, then ascending size and past-return rank; the last seven
are long and the first seven are short.

Rebalance monthly at the prior month-end. Gross risky exposure is 2.0 and net
risky exposure is zero when invested. The 100\% equity collateral earns the
risk-free rate, with no additional funding spread. Each evaluation segment starts
flat; there is no terminal liquidation.

Trading cost is 0.001 per dollar traded, and short borrowing costs 0.005 annually.
Deduct costs at month-end. With $\widetilde w_{i,t}$ denoting the drifted weight
before rebalancing,
\begin{equation}
\begin{aligned}
 r^{\ex,\mathrm{gross}}_{s,t}
    &=\sum_{i=1}^{75}w_{i,t}r_{i,t},\\
 r^{\ex,\mathrm{net}}_{s,t}
    &=r^{\ex,\mathrm{gross}}_{s,t}
      -0.001\sum_{i=1}^{75}|w_{i,t}-\widetilde w_{i,t}|
      -\frac{0.005}{12}\sum_{i:w_{i,t}<0}(-w_{i,t}),\\
 r^{\mathrm{net,total}}_{s,t}
    &=r_{f,t}+r^{\ex,\mathrm{net}}_{s,t}.
\end{aligned}
\label{eq:strategy-return}
\end{equation}
After each month, the next pretrade risky weight is
$w_{i,t}(1+r_{i,t})/(1+r^{\mathrm{net,total}}_{s,t})$.

Step 1 screened full-sample information (including the test period), so the test
period is not an independent discovery test; this rule is version 2, revised
after earlier test results were viewed.

\clearpage
\subsection*{Performance}

For each reported excess-return series, compute its arithmetic monthly mean,
sample standard deviation, and Sharpe ratio:
\begin{equation}
 \bar r^{\ex}_s=\frac{1}{N}\sum_t r^{\ex}_{s,t},\qquad
 s_s=\sqrt{\frac{1}{N-1}\sum_t(r^{\ex}_{s,t}-\bar r^{\ex}_s)^2},\qquad
 \mathrm{SR}_{m}=\frac{\bar r^{\ex}_s}{s_s},\qquad
 \mathrm{SR}_{a}=\sqrt{12}\,\mathrm{SR}_{m}.
 \label{eq:performance}
\end{equation}
Annualization uses $\sqrt{12}$ without autocorrelation adjustment.

\begin{table}[htbp]
\centering
\caption{Strategy performance using the selected rule}
\label{tab:performance}
\small
\begin{tabular}{llrrrr}
\toprule
Sample & Excess returns & Mean (\%) & SD (\%) & Monthly SR & Annualized SR\\
\midrule
Training walk-forward & Gross & 1.434 & 3.237 & 0.443 & 1.534\\
                      & Net   & 1.243 & 3.221 & 0.386 & 1.336\\
\addlinespace
Training final-frozen & Gross & 1.474 & 3.271 & 0.450 & 1.560\\
                      & Net   & 1.272 & 3.260 & 0.390 & 1.352\\
\addlinespace
Test frozen           & Gross & 0.382 & 3.837 & 0.100 & 0.345\\
                      & Net   & 0.120 & 3.835 & 0.031 & 0.108\\
\bottomrule
\end{tabular}
\par\smallskip
\begin{minipage}{\textwidth}
\footnotesize
Means and standard deviations are monthly. Training performance covers
March 1953--June 2001 ($N=580$); testing covers July 2001--June 2026 ($N=300$).
Training walk-forward re-estimates every 12 months using prior data only.
Training final-frozen applies the final training coefficients to those same
training months and is in-sample, not an unbiased backtest. Test frozen keeps
the final training coefficients fixed. Net returns subtract trading and
borrowing costs specified above.
\end{minipage}
\end{table}

Net annualized Sharpe: training walk-forward 1.336; final frozen rule over the
training months 1.352 (in-sample, not an unbiased backtest); test 0.108; test
performance is positive but weaker (HAC mean $t=0.61$).
The HAC mean $t$-statistic uses six monthly lags.

\end{document}
